\documentclass[aspectratio=169,10pt]{beamer}
\usepackage[UTF8,fontset=fandol]{ctex}
\usepackage{amsmath,amssymb,booktabs,tabularx,array,tikz,pgfplots,pgfpages}
\usetikzlibrary{arrows.meta,positioning}
\pgfplotsset{compat=1.18}
\linespread{1.12}
\renewcommand{\arraystretch}{1.15}

\definecolor{accent}{HTML}{28649A}
\definecolor{ink}{HTML}{17334F}
\setbeamercolor{normal text}{fg=ink,bg=white}
\setbeamercolor{structure}{fg=accent}
\setbeamercolor{frametitle}{fg=ink}
\setbeamercolor{item}{fg=accent}
\setbeamerfont{frametitle}{size=\fontsize{13.5}{17}\selectfont,series=\bfseries}
\setbeamerfont{title}{size=\Huge,series=\bfseries}
\setbeamersize{text margin left=9mm,text margin right=9mm}
\setbeamertemplate{navigation symbols}{}
\setbeamertemplate{itemize items}[circle]
\setbeamertemplate{frametitle}{\vspace{3mm}\insertframetitle\par\vspace{1mm}\textcolor{accent}{\rule{\textwidth}{0.5pt}}}

\newcommand{\currentpresenter}{第 8 队}
\newcommand{\pageindicator}{\insertframenumber\,/\,\inserttotalframenumber}
\newcommand{\presenter}[1]{\renewcommand{\currentpresenter}{#1}}
\setbeamertemplate{footline}{%
    \vskip2mm
    \hspace{9mm}\begin{minipage}{\dimexpr\paperwidth-18mm\relax}
        \color{ink}\fontsize{7}{8}\selectfont\insertsectionhead\quad\currentpresenter
        \hfill\pageindicator
    \end{minipage}\vspace{2mm}}
\setbeamertemplate{note page}{\vspace{4mm}{\Large\bfseries\insertframetitle}\par\vspace{3mm}{\footnotesize\insertnote}}
\setbeameroption{hide notes}
% For presenter view, replace hide notes with show notes on second screen=right.

\newcolumntype{Y}{>{\raggedright\arraybackslash}X}
\newcommand{\key}[1]{\textcolor{accent}{\bfseries #1}}
\newcommand{\takeaway}[1]{\par\vspace{4mm}{\large\bfseries #1}}
\newcommand{\papersection}[3]{%
    \section{#1}\presenter{#3}
    \begin{frame}[plain]
        \vspace{7mm}{\small\color{accent}专题论文阅读}\par
        \vspace{7mm}{\fontsize{28}{34}\selectfont\bfseries #1}\par
        \vspace{4mm}{\Large #2}\par
        \vspace{7mm}\textcolor{accent}{\rule{\textwidth}{0.8pt}}\par
        \vspace{3mm}#3
        \note{计划15秒。这一部分介绍#1。#2。}
    \end{frame}}

\hypersetup{pdftitle={长上下文场景结构化记忆建模：小组学术汇报},pdfauthor={第8队}}

\hypersetup{pdftitle={长对话结构化记忆复现与初步诊断},pdfauthor={NLP Team08}}
\begin{document}
\source{项目方案；A—F 固定提交，来源见 sources.md}
\begin{frame}{长对话结构化记忆复现与初步诊断}
\small
第 8 队 · 第一次课设研讨 · 2026.10\par\vspace{6mm}
\begin{itemize}
\item 研究问题：历史越来越长，如何保留并用对事件证据？\vspace{4mm}
\item StructMem 小规模方法复现｜固定 30 道开发题\vspace{4mm}
\item 提交候选：署名、贡献比例与出场日期待全组确认\vspace{4mm}
\end{itemize}
\note{老师好，我们是第八队。我们的课设研究长对话中的结构化记忆，当前以 StructMem 为主方法。今天汇报固定三十道开发题上已经完成的系统对照，重点说明方法、真实进展，以及结果还能支持多强的结论。}
\end{frame}
\source{项目方案 §§1—4；真实案例 conv-41/20}
\begin{frame}{相关事实被找到，不代表事件被理解}
\small
\begin{columns}[T]
\begin{column}{0.48\textwidth}
\textbf{原始对话不断增长}\par\vspace{3mm}
\begin{itemize}
\item 同一人物在不同时间谈到不同事件
\item 相似词容易召回“有关、但不是这一次”的片段
\item 记录时间与事件发生时间未必相同
\end{itemize}\end{column}
\begin{column}{0.48\textwidth}
\textbf{我们的研究切入点}\par\vspace{3mm}
\begin{itemize}
\item 把事实、关系和事件来源一起保存
\item 先检索相关记忆，再让模型回答
\item 比较回答质量，也报告输入与构建成本
\end{itemize}\end{column}
\end{columns}
\aside{目标是可解释、可复跑的课程复现；不预设必须超过简单基线。}
\note{长对话问答的难点并不只是上下文放不下。一个人可能多次谈旅行、工作和健康，关键词相似的句子不一定指同一次事件。即使检索到了事实，模型也可能把说话日期当成事件日期。比如四月十日说“去年去旅行”，问旅行年份时，不能回答四月十日。我们的切入点是给记忆保留事件、人物和时间来源，再观察这种组织方式是否帮助回答，同时记录它需要增加多少输入和构建成本。}
\end{frame}
\source{A \#38；B \#37；C \#43；D \#41；F \#35/\#42/\#44}
\begin{frame}{两条回答路径，共用数据与调用入口}
\small
\begin{itemize}
\item A 固定会话与问题：答案另存，只在评分阶段关联\vspace{4mm}
\item S0：原始对话分块 → MiniLM 向量检索 top-60 → F 回答\vspace{4mm}
\item S2：事实/关系抽取 → 事件绑定 → 跨事件综合 → 60 事件 + 5 综合 → F 回答\vspace{4mm}
\item E：关联冻结参考 → 固定裁判 → 原始分数 + 独立证据复核\vspace{4mm}
\end{itemize}
\aside{S0/S2 的记忆、输入形态和方法指令不同，是系统对照；S1 尚未运行。}
\note{数据和问题首先由 A 冻结，参考答案单独保存，不进入记忆构建或回答输入。普通检索基线 S 零由 B 实现，把原始对话切块，用本地 MiniLM 编码，再检索六十个片段。完整方法 S 二先由 C 抽取事实和关系，保存说话人、时间及来源，再由 D 做跨事件综合。回答时取六十条事件记忆和五条综合记忆。两条路径都对接 F 的统一回答入口，最后由 E 关联参考答案评分。这里比较的是两套系统，不能把分差只归因于某一个模块。}
\end{frame}
\source{C 433de34；D 38214e5；F 49d8a67；E 1cfd177}
\begin{frame}{已完成真实闭环；复现差异明确保留}
\small
\begin{tabularx}{\linewidth}{@{}YYY@{}}\toprule
阶段 & 已完成的真实产物 & 口径\\[1.5mm]\midrule
A / B & 3 会话 · 30 题 · 每题 60 原始块 & 5 题为 30 题子集\\[1.5mm]
C & 91 session · 1,973 消息 → 3,712 事件记忆 & 2,309 事实 + 1,403 关系\\[1.5mm]
D & 91 条真实综合 · S2 回答 30/30 & 60 事件 + 5 综合\\[1.5mm]
F / E & S0 回答 30/30 · E 报告裁判 60/60 & 逐题裁判交接待完整复算\\[1.5mm]
\bottomrule\end{tabularx}
\aside{替换模型与嵌入；关闭压缩/话题切分；时间绑定含顺序构造。不是论文原设置逐字复现。}
\note{目前已经做完的真实工作包括：三个开发会话、九十一个会话时段，共一千九百七十三条原始消息；普通基线完成三十题向量检索。事件抽取得到三千七百一十二条记忆，之后生成九十一条综合记忆。两种系统各完成三十条真实回答，E 也报告完成六十次裁判。我们保留了相对论文的替换：使用 DeepSeek 和本地轻量嵌入，关闭压缩与话题切分，部分绑定时间表示顺序而非真实发送时刻。因此这是有明确差异的方法复现。}
\end{frame}
\source{E 1cfd177 汇总；本轮 answer-audit.json；成本未与账单核对}
\begin{frame}{开发集：S2 多答对 4 题，输入约为 4.42 倍}
\small
\begin{tabularx}{\linewidth}{@{}YYY@{}}\toprule
题型（官方编号） & S0 & S2\\[1.5mm]\midrule
1 多跳 · 8 题 & 4/8 & 4/8\\[1.5mm]
2 时间 · 8 题 & 5/8 & 7/8\\[1.5mm]
3 开放域 · 7 题 & 2/7 & 3/7\\[1.5mm]
4 单跳 · 7 题 & 4/7 & 5/7\\[1.5mm]
\bottomrule\end{tabularx}
\par\vspace{3mm}\textbf{S0：15/30（50.0\%）；S2：19/30（63.3\%）}\par 原始裁判为16/30与19/30；输入tokens约4.42倍。
\aside{E 已报告；生成记录已逐条核验，60 条裁判尚待独立复算。开发集差异不代表稳定泛化。}
\note{先看结果的口径。原始自动裁判给普通基线十六题正确，完整系统十九题正确。证据复核发现普通基线有一条日期回答自相矛盾，将它改为错误，因此展示的复核后结果是十五比十九，即百分之五十和百分之六十三点三。分类差异主要来自时间题两道，开放域和单跳各一道，多跳相同。两组报告的评分覆盖都是三十题。与此同时，完整系统回答输入约是基线的四点四二倍，回答费用估算约一元一角七分，而基线约两角七分。这些费用尚未核账。我们已经核验六十条生成记录，但 E 的逐题裁判交接仍待独立复算。这里是三十道开发题的现象，不是稳定泛化提升的证明。}
\end{frame}
\source{冻结开发历史 D11:3/D11:5、D16:8；E 六例复核；本轮源文核对}
\begin{frame}{案例帮助解释结果，也暴露评分边界}
\small
\begin{columns}[T]
\begin{column}{0.48\textwidth}
\textbf{相对时间：记录日期 $\ne$ 事件年份}\par\vspace{3mm}
\begin{itemize}
\item conv-41/20：问太平洋西北旅行时间
\item 2023-04-10 的 D11:3：“a road trip we took last year”（节选）
\item S0：10 April, 2023
\item S2：Last year (2022)
\end{itemize}\end{column}
\begin{column}{0.48\textwidth}
\textbf{错误和参考都要检查}\par\vspace{3mm}
\begin{itemize}
\item 日期矛盾题 conv-42/28：S0 改判 1 次
\item conv-42/61：来源不足以证明具体游戏平台
\item conv-41/41：海景照片不能证明住址
\item 保留参考和分母，单列证据不足
\end{itemize}\end{column}
\end{columns}
\aside{案例为 AI 辅助证据核查，人工验收未填写；观察到系统差异，不等于确定模块因果。}
\note{一个清楚的例子是公路旅行时间。原文在二零二三年四月十日说“我们去年进行的公路旅行”，下一条相关对话说明地点是太平洋西北。普通基线直接回答记录日期，完整系统回答去年，也就是二零二二年。另一道冰淇淋题在六月二十四日说“上周五”，普通基线却同时写六月二十四日和此前的星期五，因此复核修改了一次裁判标签。但并非所有错分都能归因系统。游戏平台和海滩住址两题的冻结参考缺少充分来源，我们保留原答案和分母，并单列限制。这些是人工智能辅助核查，还不是组员亲自验收。}
\end{frame}
\source{项目方案 §§3、6—7；S1 与检索预算实验均未运行}
\begin{frame}{创新性论证：从可检验的问题出发}
\small
\begin{columns}[T]
\begin{column}{0.48\textwidth}
\textbf{本轮已得到的研究依据}\par\vspace{3mm}
\begin{itemize}
\item 结构化路径在部分时间题表现不同
\item 回答输入增加约 4.42 倍，质量与成本需一起看
\item 仍存在证据不足、参考口径和判分可靠性问题
\end{itemize}\end{column}
\begin{column}{0.48\textwidth}
\textbf{研讨后的验证顺序（未执行）}\par\vspace{3mm}
\begin{itemize}
\item 先补 S1：同事件库、同指令，关闭跨事件综合
\item 再固定比较检索 30 与 60，综合数保持 5
\item 按正确性、有效覆盖与输入成本决定是否保留
\end{itemize}\end{column}
\end{columns}
\aside{课程贡献是可复跑的适配、对照与证据分析；尚不声称提出学术新方法。}
\note{创新性不能只靠给系统起一个新名字。现阶段我们的贡献是完成可复跑的适配和系统对照，并发现一个具体问题：结构化方法多用输入，是否值得，收益来自哪里？下一步先补事件记忆但关闭跨事件综合的 S 一，保持事件库和方法指令一致，检查综合模块的作用。之后才考虑把检索数量从六十减到三十，综合数量保持五条，比较质量和输入成本。两项都尚未执行，也不会在看见少量题变好后，就宣布方法具有稳定优势。}
\end{frame}
\source{项目方案 §8；Issues \#19/\#24/\#26/\#28/\#31/\#32；成员原稿}
\begin{frame}{六人已交付内容与剩余协作事项}
\small
{\footnotesize\setlength{\tabcolsep}{2pt}\renewcommand{\arraystretch}{1.02}
\begin{tabularx}{\linewidth}{@{}p{43mm}Yp{22mm}@{}}\toprule
任务 / 账号 & 实际贡献 & 当前比例\\[0.3mm]\midrule
A · ngocnamd93-spec & 冻结题单、答案隔离、嵌入环境 & 待全组确认\\[0.3mm]
B · YaoYang-123456 & 原始块 top-60、真实检索、MetaMem 阅读 & 待全组确认\\[0.3mm]
C · oioioi-521 & 事实/关系抽取、事件绑定与持久化 & 待全组确认\\[0.3mm]
D · lanshuheng9-oss & 真实综合、S2 检索与回答 & 待全组确认\\[0.3mm]
E · onlysonder & 固定裁判、结果表与案例复核材料 & 待全组确认\\[0.3mm]
F · hycx233 & 统一客户端、S0 回答、交接及全组整合 & 待全组确认\\[0.3mm]
\bottomrule\end{tabularx}
}
\aside{贡献比例不按发言时长或提交数推算。主讲、录制、上传负责人尚未确认。}
\note{六位成员的实际分工在表中列出。A 负责数据冻结和环境；B 负责原始对话检索及个人 MetaMem 阅读；C 构建事件记忆；D 做综合、检索和完整系统回答；E 固定评分并整理案例；F 提供统一调用、完成基线回答并负责材料整合。我们没有按提交数量或发言时间自动分摊贡献比例，这一栏需要全组根据实际工作确认。当前风险也很具体：评分程序的输入绑定和中断保存需要补丁审核，逐题裁判交接需要补全，真人排练与视频负责人尚未确认。这些交付事项与继续调模型是两回事。}
\end{frame}
\source{教师通知：22/26/29 日；云盘有效期 22 日 14:30（北京时间）}
\begin{frame}{现在可以解释成果；下一步补齐验收与消融}
\small
\begin{itemize}
\item 已完成：固定数据 → 两套真实回答 → 自动评分报告与证据案例\vspace{4mm}
\item 当前观察：S2 在本开发集多答对 4 题，代价是更大的输入\vspace{4mm}
\item 提交前：逐题评分交接、成员署名与比例、排练及视频安排\vspace{4mm}
\item 研讨后：先做 S1 消融，再决定是否开展检索预算实验\vspace{4mm}
\end{itemize}
\aside{按 10 月 21 日前可展示准备；实际场次以课程 Excel 为准。视频由组员提交。}
\note{总结来说，我们已经建立固定数据到两套真实回答的闭环，并有自动评分报告和可追溯案例。当前观察是完整系统在开发集上多答对四题，但输入明显增加，因此不能只报一个更高分数。提交前优先补齐评分交接、署名比例和排练录制安排，不再为成绩好看增加实验。研讨之后先完成 S 一消融，再决定检索预算实验是否值得做。我们会据证据说明已经完成什么、还不能证明什么。谢谢老师。}
\end{frame}
\appendix
\source{C/D/F/E 固定汇总；失败用量未知不按 0 处理}
\begin{frame}{备用：费用、计数与评价版本}
\small
\begin{tabularx}{\linewidth}{@{}YYY@{}}\toprule
阶段 & 请求 / token & 已记录费用估算\\[1.5mm]\midrule
C 事件抽取 & 182 成功；另 2 次连接失败、用量未知 & 单列，未纳入下方小计\\[1.5mm]
D 跨事件综合 & 91 次；339,859 tokens & ¥0.894722\\[1.5mm]
F / D 回答 & S0 30 次；S2 30 次 & ¥0.268430 / ¥1.173988\\[1.5mm]
E 自动裁判 & 60 次；9,353 tokens & 未核账，未纳入小计\\[1.5mm]
回答 + 综合小计 & 不包含 C 抽取和 E 裁判 & ¥2.337140\\[1.5mm]
\bottomrule\end{tabularx}
\aside{不是项目总账单。S0 16→15 是一条复核修订，不能混写原始与修订成绩。}
\note{}
\end{frame}
\source{项目方案、教师通知、\#46 审核意见}
\begin{frame}{备用：对照边界与仍待确认的事项}
\small
\begin{itemize}
\item S0：原始块向量检索；S1：仅事件记忆；S2：事件 + 跨事件综合\vspace{4mm}
\item S1 未运行；S0/S2 还存在方法提示词与上下文体量差异\vspace{4mm}
\item 原始裁判 $\ne$ AI 证据复核 $\ne$ 成员人工验收\vspace{4mm}
\item 尚待：逐题裁判交接、真实署名/比例、Excel 顺序、场次、格式与视频负责人\vspace{4mm}
\end{itemize}
\aside{8 分钟课设研讨视频与 11 月 1 日约 60 分钟阅读视频是不同交付。}
\note{}
\end{frame}
\end{document}
